% GENERATED FILE. Edit canonical lessons, then run npm run generate:books.
% canonical-source-hash: fnv1a64:6920e9d8239fce16
% canonical-lessons: LA-C383-magistratus, LA-C383-praetor, LA-C383-aedilis, LA-C383-quaestor, LA-C383-censor

\chapter{The Magistrates}
\label{ch:la-the-magistrates}
\hlchaptermodality{383}

\begin{chapteropening}
\textbf{By the end of this chapter:} \emph{I can name a magistrate, a praetor, an aedile, a quaestor and a censor, and say what each one does.}

Complete the last lesson of chapter 383: I can name a magistrate, a praetor, an aedile, a quaestor and a censor, and say what each one does.
\end{chapteropening}

\section[magistrātus, magistrātūs]{magistrātus, magistrātūs — a public office; a magistrate}

\label{lesson:LA-C383-magistratus}



\begin{quote}
Before the new one: say the Latin for well known, then the Latin for of the people.
\end{quote}

\subsection*{You'll want to know: magistrātus, magistrātūs}
\textbf{magistrātus, magistrātūs} — \textquotedblleft{}a public office; a magistrate\textquotedblright{}.

\begin{quote}
Magistrātūs quotannīs ā populō creantur.
\end{quote}

\begin{quote}
\emph{Magistrates are elected by the people every year.}
\end{quote}

\emph{magistrātus, magistrātūs} m., built on \emph{magister}: both the office and the man who holds it. Roman offices lasted a year.

\begin{cousinweb}
\emph{magistrātus} is built on \emph{magister} \textquotedblleft{}a master\textquotedblright{} (EtymDB; Lewis \& Short): the office of one set over others. English \emph{magistrate} comes from it.
\end{cousinweb}

\begin{grammarlens}[title={What you've built}]
One more word for telling, explaining and describing at B1.
\end{grammarlens}

\subsection*{Your turn}
\begin{itemize}
\raggedright
  \item \emph{Say it:} \emph{magistrātus, magistrātūs}
  \item \emph{Say it:} \emph{magistrātus, magistrātūs}, once more
  \item \emph{Say it:} a sentence of your own with \emph{magistrātus, magistrātūs}
  \item \emph{Recall:} say the Latin for well known, then the Latin for of the people, then say \emph{magistrātus, magistrātūs} again
\end{itemize}

\begin{tcolorbox}[breakable,colback=teal!4,colframe=teal!35!black,title={Before you move on}]
What does magistrātus, magistrātūs mean? (\textquotedblleft{}A public office; a magistrate\textquotedblright{}.) Say it once more.
\end{tcolorbox}
\section[praetor, praetōris]{praetor, praetōris — a praetor, a judge and magistrate}

\label{lesson:LA-C383-praetor}



\begin{quote}
Before the new one: say the Latin for of the people, then the Latin for a public office.
\end{quote}

\subsection*{You'll want to know: praetor, praetōris}
\textbf{praetor, praetōris} — \textquotedblleft{}a praetor, a judge and magistrate\textquotedblright{}.

\begin{quote}
Praetor iūs dīcit in forō.
\end{quote}

\begin{quote}
\emph{The praetor gives judgement in the Forum.}
\end{quote}

\emph{praetor, praetōris} m. Next in rank below the consuls, he ran the law courts.

\begin{cousinweb}
\emph{praetor} is usually taken as \emph{prae} \textquotedblleft{}before\textquotedblright{} and the root of \emph{eō} \textquotedblleft{}to go\textquotedblright{} (Lewis \& Short): one who goes in front, an old word for a leader.
\end{cousinweb}

\begin{grammarlens}[title={What you've built}]
One more word for telling, explaining and describing at B1.
\end{grammarlens}

\subsection*{Your turn}
\begin{itemize}
\raggedright
  \item \emph{Say it:} \emph{praetor, praetōris}
  \item \emph{Say it:} \emph{praetor, praetōris}, once more
  \item \emph{Say it:} a sentence of your own with \emph{praetor, praetōris}
  \item \emph{Recall:} say the Latin for of the people, then the Latin for a public office, then say \emph{praetor, praetōris} again
\end{itemize}

\begin{tcolorbox}[breakable,colback=teal!4,colframe=teal!35!black,title={Before you move on}]
What does praetor, praetōris mean? (\textquotedblleft{}A praetor, a judge and magistrate\textquotedblright{}.) Say it once more.
\end{tcolorbox}
\section[aedīlis, aedīlis]{aedīlis, aedīlis — an aedile, an officer for streets, markets and games}

\label{lesson:LA-C383-aedilis}



\begin{quote}
Before the new one: say the Latin for a public office, then the Latin for a praetor.
\end{quote}

\subsection*{You'll want to know: aedīlis, aedīlis}
\textbf{aedīlis, aedīlis} — \textquotedblleft{}an aedile, an officer for streets, markets and games\textquotedblright{}.

\begin{quote}
Aedīlis macellum et viās cūrat.
\end{quote}

\begin{quote}
\emph{The aedile looks after the meat market and the streets.}
\end{quote}

\emph{aedīlis, aedīlis} m., built on the word for a temple or a building: he first cared for the temples. English \emph{aedile}.

\begin{cousinweb}
\emph{aedīlis} is built on \emph{aedēs} \textquotedblleft{}a temple, a building\textquotedblright{} (EtymDB, which lists old Latin \emph{aidīlis}; Lewis \& Short). English \emph{aedile} comes from it.
\end{cousinweb}

\begin{grammarlens}[title={What you've built}]
One more word for telling, explaining and describing at B1.
\end{grammarlens}

\subsection*{Your turn}
\begin{itemize}
\raggedright
  \item \emph{Say it:} \emph{aedīlis, aedīlis}
  \item \emph{Say it:} \emph{aedīlis, aedīlis}, once more
  \item \emph{Say it:} a sentence of your own with \emph{aedīlis, aedīlis}
  \item \emph{Recall:} say the Latin for a public office, then the Latin for a praetor, then say \emph{aedīlis, aedīlis} again
\end{itemize}

\begin{tcolorbox}[breakable,colback=teal!4,colframe=teal!35!black,title={Before you move on}]
What does aedīlis, aedīlis mean? (\textquotedblleft{}An aedile, an officer for streets, markets and games\textquotedblright{}.) Say it once more.
\end{tcolorbox}
\section[quaestor, quaestōris]{quaestor, quaestōris — a quaestor, a finance officer}

\label{lesson:LA-C383-quaestor}



\begin{quote}
Before the new one: say the Latin for a praetor, then the Latin for an aedile.
\end{quote}

\subsection*{You'll want to know: quaestor, quaestōris}
\textbf{quaestor, quaestōris} — \textquotedblleft{}a quaestor, a finance officer\textquotedblright{}.

\begin{quote}
Quaestor pecūniam pūblicam servat.
\end{quote}

\begin{quote}
\emph{The quaestor keeps the public money.}
\end{quote}

\emph{quaestor, quaestōris} m., built on \emph{quaerō}, to seek, to inquire. The lowest regular office: the treasury and the accounts.

\begin{cousinweb}
\emph{quaestor} is built on \emph{quaerō} \textquotedblleft{}to seek, to inquire\textquotedblright{} (Lewis \& Short), with the ending \emph{-tor} of the one who does it.
\end{cousinweb}

\begin{grammarlens}[title={What you've built}]
One more word for telling, explaining and describing at B1.
\end{grammarlens}

\subsection*{Your turn}
\begin{itemize}
\raggedright
  \item \emph{Say it:} \emph{quaestor, quaestōris}
  \item \emph{Say it:} \emph{quaestor, quaestōris}, once more
  \item \emph{Say it:} a sentence of your own with \emph{quaestor, quaestōris}
  \item \emph{Recall:} say the Latin for a praetor, then the Latin for an aedile, then say \emph{quaestor, quaestōris} again
\end{itemize}

\begin{tcolorbox}[breakable,colback=teal!4,colframe=teal!35!black,title={Before you move on}]
What does quaestor, quaestōris mean? (\textquotedblleft{}A quaestor, a finance officer\textquotedblright{}.) Say it once more.
\end{tcolorbox}
\section[cēnsor, cēnsōris]{cēnsor, cēnsōris — a censor}

\label{lesson:LA-C383-censor}



\begin{quote}
Before the new one: say the Latin for an aedile, then the Latin for a quaestor.
\end{quote}

\subsection*{You'll want to know: cēnsor, cēnsōris}
\textbf{cēnsor, cēnsōris} — \textquotedblleft{}a censor\textquotedblright{}.

\begin{quote}
Cēnsor cīvēs et bona eōrum numerat.
\end{quote}

\begin{quote}
\emph{The censor counts the citizens and their property.}
\end{quote}

\emph{cēnsor, cēnsōris} m., built on \emph{cēnseō}, to assess. Every five years two censors drew up the list of citizens and could strike a senator off the roll. English \emph{censor}.

\begin{cousinweb}
\emph{cēnsor} is built on \emph{cēnseō} \textquotedblleft{}to assess, to value\textquotedblright{} (EtymDB; Lewis \& Short): the officer who assessed the citizens. English \emph{censor} comes from it.
\end{cousinweb}

\begin{grammarlens}[title={What you've built}]
One more word for telling, explaining and describing at B1.
\end{grammarlens}

\subsection*{Your turn}
\begin{itemize}
\raggedright
  \item \emph{Say it:} \emph{cēnsor, cēnsōris}
  \item \emph{Say it:} \emph{cēnsor, cēnsōris}, once more
  \item \emph{Say it:} a sentence of your own with \emph{cēnsor, cēnsōris}
  \item \emph{Recall:} say the Latin for an aedile, then the Latin for a quaestor, then say \emph{cēnsor, cēnsōris} again
\end{itemize}

\begin{tcolorbox}[breakable,colback=teal!4,colframe=teal!35!black,title={Before you move on}]
What does cēnsor, cēnsōris mean? (\textquotedblleft{}A censor\textquotedblright{}.) Say it once more.
\end{tcolorbox}
